\documentclass[11pt]{article}

\usepackage{amsmath,amssymb,amsthm,mathtools}
\usepackage[T1]{fontenc}
\usepackage{lmodern}
\usepackage[margin=1in]{geometry}
\usepackage{enumitem}
\usepackage{hyperref}
\usepackage{microtype}

\hypersetup{
  colorlinks=true,
  linkcolor=blue,
  citecolor=blue,
  urlcolor=blue
}

\newtheorem{theorem}{Theorem}[section]
\newtheorem{proposition}[theorem]{Proposition}
\newtheorem{corollary}[theorem]{Corollary}
\newtheorem{definition}[theorem]{Definition}
\newtheorem{lemma}[theorem]{Lemma}
\newtheorem{remark}[theorem]{Remark}

\newcommand{\Z}{\mathbb Z}

\title{A compact reproof of the 079 chart/interface landing}
\author{}
\date{March 2026}

\begin{document}
\maketitle

\begin{abstract}
This note removes the standalone 079 support note by deriving the exact
chart/interface statement used in 092 from the already frozen regular-slice
chart formulas.

The proof stays at the correct scope. It does \emph{not} prove the actual
post-exit continuation exists; that still comes from the structural first-exit
package (062, written explicitly in 092 Theorem~3.1). What this note does
prove is that once that post-exit continuation is read in the grouped
full-chain chart, its first two labels are forced to be
\[
3m-2 \to 3m-1.
\]
Equivalently, the exceptional chain-label continuation is
\[
3m-3 \to 3m-2 \to 3m-1.
\]
So after 095, the cleaned 092 suite no longer needs 079 as a separate
accepted-support theorem.
\end{abstract}

\section{Scope and imported chart package}

Fix an odd modulus $m$ in the accepted D5 regime and the best-seed channel
\[
R1 \to H_{L1}.
\]

The present note uses only the already frozen regular-slice chart package:
\begin{enumerate}[label=\textup{(\arabic*)},leftmargin=2.5em]
\item for each regular source label $u$, the $\Theta=2$ window starts at
      \[
      s_u = m((u+2) \bmod m)-1;
      \]
\item each such $\Theta=2$ window has exactly $m(m-3)$ consecutive full-chain
      labels;
\item grouped successor on the chart quotient is exact: whenever the
      continuation is read in that quotient, labels advance by
      \[
      \delta \mapsto \delta+1.
      \]
\end{enumerate}

These formulas are the stable chart-side data already frozen in
\begin{itemize}[leftmargin=2em]
\item \texttt{checks/d5\_077\_compact\_interval\_summary.json}, and
\item \texttt{checks/d5\_081\_regular\_union\_endpoint\_table.csv}.
\end{itemize}

The present note uses the structural first-exit theorem only at the very end,
when connecting the chart labels back to the actual exceptional post-exit
continuation.

\section{Algebra of the regular source windows}

\begin{definition}[start and terminal labels]
For a regular source label $u$, write
\[
s_u = m((u+2) \bmod m)-1
\]
for the $\Theta=2$ start label of the source-$u$ window.

Because that window has length $m(m-3)$, its terminal label is
\[
e_u = s_u + m(m-3)-1
\]
in $\Z/m^2\Z$.
\end{definition}

\begin{proposition}[closed formula for the terminal label]
For every regular source label $u$,
\[
e_u = m(u-1)-2 \pmod{m^2}.
\]
\end{proposition}

\begin{proof}
Substitute the start formula into the terminal formula:
\[
e_u = m((u+2) \bmod m)-1 + m(m-3)-1.
\]
Modulo $m^2$, this is
\[
e_u = m(u-1)-2.
\]
That is exactly the endpoint formula already recorded in the promoted regular
endpoint table.
\end{proof}

\begin{corollary}[the next source window is shifted by $-3$]
The label immediately after the terminal source-$u$ label is the start of the
source-$u-3$ window:
\[
e_u + 1 = s_{u-3}
\]
(modulo $m$, with representatives in $\{1,2,4,5,\dots,m-1\}$).
\end{corollary}

\begin{proof}
From Proposition~2.2,
\[
e_u + 1 = m(u-1)-1.
\]
But
\[
s_{u-3} = m(((u-3)+2) \bmod m)-1 = m(u-1)-1.
\]
So the two labels agree in $\Z/m^2\Z$.
\end{proof}

\begin{corollary}[the common regular interface]
The terminal regular source-$4$ label is
\[
e_4 = 3m-2,
\]
and the regular source-$1$ start label is
\[
s_1 = 3m-1.
\]
Hence the common regular interface is exactly
\[
3m-2 \to 3m-1.
\]
\end{corollary}

\begin{proof}
From Proposition~2.2,
\[
e_4 = m(4-1)-2 = 3m-2.
\]
From the start formula,
\[
s_1 = m(1+2)-1 = 3m-1.
\]
By Corollary~2.3, the label after $e_4$ is $s_1$. So the interface is exactly
\[
3m-2 \to 3m-1.
\]
Because each regular source window is one contiguous interval, $e_4$ is the
unique terminal source-$4$ representative and $s_1$ is the source-$1$ start
representative.
\end{proof}

\section{The reproof replacing 079}

\begin{theorem}[chart/interface landing]
Read the actual exceptional post-exit continuation in the grouped full-chain
chart used for the boundary labels. Then its first two chart labels are forced
to be
\[
3m-2 \to 3m-1.
\]
Equivalently, the exceptional chain-label continuation is
\[
3m-3 \to 3m-2 \to 3m-1,
\]
where $3m-2$ is the unique terminal regular source-$4$ occurrence and $3m-1$
is the regular source-$1$ start.
\end{theorem}

\begin{proof}
The exceptional cutoff row is, by definition,
\[
\delta = 3m-3.
\]
In the grouped chart quotient, successor is exact and advances the label by
$+1$. Therefore the first chart label after the cutoff row is forced to be
\[
3m-2.
\]
By Corollary~2.4, that label is the common regular interface point: the unique
terminal regular source-$4$ representative, followed by the regular source-$1$
start at
\[
3m-1.
\]
So the chain-label continuation is forced to be
\[
3m-3 \to 3m-2 \to 3m-1.
\]
This is exactly the chart/interface statement that 092 uses later.
\end{proof}

\begin{corollary}[how this combines with 062]
Combine Theorem~3.1 with the structural first-exit theorem of 062 / 092:
\begin{itemize}[leftmargin=2em]
\item 062 supplies existence of the actual exceptional post-exit continuation
      and its universal raw exit target;
\item Theorem~3.1 identifies the corresponding full-chain labels as
      \[
      3m-3 \to 3m-2 \to 3m-1.
      \]
\end{itemize}
So the old 062+079 composition step becomes 062+095, with the same honest
scope: raw existence comes from 062, and chart/interface label identification
comes from the present note.
\end{corollary}

\section{Exact finite support table}

The frozen chart support files record the source-$4$ terminal label and the
source-$1$ start label exactly at the interface values below.

\begin{center}
\begin{tabular}{rccc}
\hline
$m$ & $3m-3$ & source-$4$ terminal $=3m-2$ & source-$1$ start $=3m-1$ \\
\hline
5  & 12 & 13 & 14 \\
7  & 18 & 19 & 20 \\
9  & 24 & 25 & 26 \\
11 & 30 & 31 & 32 \\
13 & 36 & 37 & 38 \\
15 & 42 & 43 & 44 \\
17 & 48 & 49 & 50 \\
19 & 54 & 55 & 56 \\
21 & 60 & 61 & 62 \\
\hline
\end{tabular}
\end{center}

For the original 079 support range $m=13,15,17,19,21$, this matches the
promoted chart statement exactly. It also agrees with the frozen
smaller-modulus regular-slice data, including the exact $m=5$ interface
\[
12 \to 13 \to 14.
\]

\section{What changes in the cleaned package}

\begin{corollary}[accepted-support list after 095]
After inserting Theorem~3.1 into the cleaned 092 suite, the standalone 079
interface note is no longer needed as a separate accepted-support theorem.

The remaining imported support in the cleaned odd-$m$ D5 package is:
\begin{enumerate}[label=\textup{(\arabic*)},leftmargin=2.5em]
\item 077 fixed-$\delta$ tail-length reduction;
\item 081 regular continuation / regular-union theorem;
\item the localized structural 033/062 input already used in the first-exit
      block.
\end{enumerate}

The chart formulas used in the present note are already frozen in the promoted
regular-slice support files, so the remaining stabilization work no longer
needs a separate 079 theorem package.
\end{corollary}

\section{Bottom line}

095 does not prove the final raw gluing theorem by itself. What it does prove
is exactly the chain-label statement that 079 was carrying:
\begin{itemize}[leftmargin=2em]
\item source-$4$ terminal regular interface label is $3m-2$;
\item source-$1$ regular start label is $3m-1$;
\item therefore the exceptional chain-label continuation is forced to pass
      through
      \[
      3m-3 \to 3m-2 \to 3m-1.
      \]
\end{itemize}
That is the only role 079 needed to play inside the cleaned 092 suite.

\end{document}
